\section{Target users}
\label{sec:users}

\subsection{Who this is for}

Learning Mocha is built for the \textbf{self-directed technical learner}: a student or working developer
who reads documentation, courses, papers and videos on their own initiative and wants the result to
accumulate into something they can navigate a year later. Concretely, the app assumes a user who

\begin{itemize}
  \item studies subjects with real internal structure --- programming, systems, mathematics, languages
        --- where topic B genuinely depends on topic A;
  \item writes their own notes rather than collecting other people's, and is comfortable with Markdown;
  \item cares that their notes stay theirs: no account, no vendor, no silent upload;
  \item reads on a phone in short sessions, so ``where was I'' and ``what should I read next'' matter more
        than word-processing features.
\end{itemize}

\subsection{Three concrete users}

\begin{mochabox}[title={Minh --- second-year CS student}]
Studying for a distributed-systems course. He types notes during lectures and wants them to end up in the
right place in a tree he can revise from. He uses \textbf{status} to mark which lecture notes are still
half-understood, the \textbf{glossary} because half the vocabulary is new, and he asks the assistant
``based on what I already have, what am I missing before Raft?'' The answer arrives as a proposed tree of
posts he can prune before accepting.
\end{mochabox}

\begin{mochabox}[title={Lan --- backend developer changing stack}]
Moving from Node to Spring Boot. She keeps a \emph{Backend} branch and writes a post per concept as she
hits it, linking \code{[[Dependency Injection]]} from five different articles. \textbf{Backlinks} tell her
which of her own notes depend on a concept she is about to revise; the \textbf{knowledge graph} shows her
that her \emph{Java} cluster is dense and her \emph{Deployment} cluster is three isolated dots. She uses
\textbf{Organize} mode to ask the assistant to tidy a branch that grew by accretion.
\end{mochabox}

\begin{mochabox}[title={Hoa --- learning a language and reading papers}]
Reads English-language material and stores Vietnamese meanings alongside definitions, which is why the
glossary has a dedicated \code{meaningVi} field rather than one free-text definition. She attaches
\textbf{YouTube resources} to posts and watches them inside the article instead of losing her place to the
browser, and she \textbf{exports} her library to Drive every week because it is the only copy.
\end{mochabox}

\subsection{Who this is not for}

Saying no is part of the design. Learning Mocha is not a team wiki (no sharing, no permissions, no
conflict resolution), not a general note-taking app (no images, no attachments, no handwriting, no
web clipper), and not a flashcard app (no spaced repetition --- \code{status} is a reading state, not a
memory schedule). Each of those would have pulled the architecture towards accounts and a server, which
is exactly what commitment~1 in \S\ref{sec:intro} rules out.

\subsection{What the user has to do, and what the app does for them}

\begin{tabularx}{\linewidth}{L{5.4cm} X}
\toprule
\textbf{The user does this} & \textbf{So the app can do this} \\
\midrule
Writes \code{[[Spring Boot]]} inside a post &
Resolves it to a real post, renders it as a tap target, records the reverse edge, and keeps it correct
when either post is renamed \\
\addlinespace
Sets a status and a few tags &
Powers ``continue reading'' on the home screen, the related-posts list, tag browsing, and search filters \\
\addlinespace
Pastes a YouTube URL into the body &
Detects it and offers it as a playable resource card on the article \\
\addlinespace
Adds a glossary term &
Surfaces it as a chip on every article that mentions it, and makes it findable from search \\
\addlinespace
Asks the assistant for a learning path &
Reads only the parts of the library it explicitly queries, proposes a reviewable tree of posts with
content, links, tags and glossary entries, and writes nothing until it is approved \\
\bottomrule
\end{tabularx}
